% Generated draft: requires mathematical and source review.
\documentclass[11pt]{article}
\usepackage[margin=25mm]{geometry}
\usepackage{fontspec}
\setmainfont{Latin Modern Roman}
\usepackage{amsmath,amssymb,mathtools,unicode-math}
\setmathfont{Latin Modern Math}
\usepackage{enumitem,needspace,xurl,hyperref}
\setlength{\parindent}{0pt}
\setlength{\parskip}{5pt}
\setlength{\emergencystretch}{3em}
\newcommand{\N}{\mathbb{N}}
\newcommand{\Z}{\mathbb{Z}}
\newcommand{\Q}{\mathbb{Q}}
\newcommand{\R}{\mathbb{R}}
\newcommand{\C}{\mathbb{C}}
\newcommand{\F}{\mathbb{F}}
\newcommand{\A}{\mathbb{A}}
\newcommand{\PP}{\mathbb{P}}
\newcommand{\E}{\mathbb{E}}
\newcommand{\eps}{\varepsilon}
\newcommand{\dd}{\,\mathrm{d}}
\begin{document}
\section*{UnsolvedMath ID: 2302006}
\textbf{Source problem code:} catalogue entry
\subsection{Definitions and mathematical statement}
For an entire function $f$, define $M(r,f)=\max_{|z|=r}|f(z)|$. A path to infinity means a continuous map $\gamma:[0,\infty)\to\C$ with $|\gamma(t)|\to\infty$. The source attributes to Boas the existence of one such path satisfying
\[
\forall n\in\Z_{\ge0},\qquad
\lim_{t\to\infty}\frac{|f(\gamma(t))|}{|\gamma(t)|^n}=\infty.
\]
The same path must work for every fixed $n$; this is not a choice of a different path for each power.

The extract says only that $f$ is entire. This literal scope includes polynomials, for which the displayed growth beyond every polynomial power cannot hold. Thus a transcendental-entire restriction, or another exclusion of polynomial functions, is missing from the stated base result. That necessary qualification is flagged explicitly rather than silently inserted into the attributed theorem. The refinement question is naturally read in the intended nonpolynomial setting, pending source confirmation.

The open request is whether quantitative information on the lower growth of $M(r,f)$ can strengthen the growth guaranteed on such a path. The extract specifies neither a class of lower-growth bounds nor a particular proposed stronger conclusion. For context, ordinary order and lower order in the transcendental setting are
\[
\rho(f)=\limsup_{r\to\infty}\frac{\log\log M(r,f)}{\log r},\qquad
\lambda(f)=\liminf_{r\to\infty}\frac{\log\log M(r,f)}{\log r}.
\]
But ``lower growth'' need not mean only the single parameter $\lambda(f)$; it could prescribe a more detailed lower bound. No such interpretation or optimal rate is imposed here.

The source also attributes to Hayman entire functions of infinite order, described as growing arbitrarily rapidly, such that every path $\gamma$ to infinity on which $|f(\gamma(t))|\to\infty$ satisfies
\[
\log\log|f(\gamma(t))|=O(\log|\gamma(t)|)\qquad(t\to\infty).
\]
The implicit bound may depend on the path; the source gives no uniform constant. This is its stated obstruction to transferring rapid maximum-modulus growth directly to a path. The phrase ``arbitrarily rapidly'' is not quantified, so no eventual pointwise lower bound against every prescribed growth function is invented. The supplied summary calls the Boas existence result classical and the quantitative refinement open, but no proof or independent source verification is supplied here.
\subsection{Short English statement}
Can lower bounds for an entire function's maximum modulus force faster growth along one path than growth beyond every fixed polynomial power? The base path statement needs a missing nonpolynomial hypothesis. The source also reports very rapidly growing infinite-order examples whose growth along every path tending to infinity in value still has finite order, so the exact lower-growth assumption and desired improvement matter.
\subsection{Sources}
\begin{itemize}
\item Local source registry: \url{data/unsolvedmath/problems.json}, permanent UnsolvedMath ID $2302006$, source code catalogue entry, including the supplied PARTIAL-PROGRESS summary and its distinction between a base theorem and an open refinement. Dataset attribution: Ulam AI and UnsolvedMath Contributors; original-source attribution and rights are retained.
\item Hayman, \emph{Research Problems in Function Theory}, source-list date 2018, Problem 2.6: \url{https://arxiv.org/abs/1809.07200}. The supplied metadata reports source-TeX extraction and July 29, 2026 access, with current-status and rights review flags.
\item Boas's path-existence result is explicitly labeled unpublished in the statement. Hayman's infinite-order examples are also asserted there without a separate complete citation. The later summary calls the base result established but provides no publication or proof resolving the missing hypotheses. No linked source was inspected for this draft.
\end{itemize}
\end{document}
